\section*{8. Recolección y análisis de datos}
 
Este capítulo detalla el proceso crítico de transformación de los datos en bruto
obtenidos mediante técnicas de extracción web hasta la consolidación de un dataset
estructurado apto para el entrenamiento de modelos de Machine Learning. Se describe
la procedencia de los datos, el proceso de integración entre fuentes heterogéneas, la
limpieza, y la ingeniería de variables que permite al sistema capturar la naturaleza
estocástica del fútbol.
 
\subsection{8.1 Fuentes de datos y proceso de scraping}
 
La construcción del dataset se fundamenta en la extracción automatizada de datos desde
tres fuentes especializadas y complementarias:
 
\begin{itemize}
    \item \textbf{FBRef.com}: Fuente principal de métricas avanzadas de rendimiento
    individual (xG, PSxG, conducciones progresivas, intercepciones, estadísticas
    defensivas y de portería). Los datos se obtienen procesando los archivos HTML de
    cada partido de las últimas 2,5 temporadas (2023/24, 2024/25 y la actual 2025/26).
 
    \item \textbf{FutbolFantasy.com}: Proporciona el contexto específico del juego
    Fantasy, incluyendo los puntos obtenidos por cada jugador en cada jornada. Constituye
    la fuente de la variable objetivo (\texttt{target\_pf\_next}).
 
    \item \textbf{Transfermarkt}: Fuente de contexto competitivo, empleada para extraer
    las rachas de los equipos, la clasificación por jornada, datos de estadios y
    nacionalidades de los jugadores.
\end{itemize}
 
\subsubsection{8.1.1 Barreras técnicas y estrategias de mitigación}
 
Durante el proceso de extracción se detectaron dos barreras críticas que obligaron a
diseñar estrategias específicas:
 
\begin{itemize}
    \item \textbf{Bloqueos de IP en FBRef}: Este portal implementa políticas extremadamente
    estrictas de \textit{rate limiting}. Ante la imposibilidad de realizar un scraping masivo
    automatizado sin ser bloqueado, se optó por la descarga controlada y local de los archivos
    HTML de cada partido, procesándolos de forma diferida. Esta decisión garantizó la
    integridad de los datos sin comprometer la disponibilidad del sistema.
 
    \item \textbf{Evasión de Cloudflare en FutbolFantasy}: El portal bloquea agentes no
    humanos. Se implementó el módulo \texttt{descargar\_puntos.py} utilizando la librería
    \texttt{cloudscraper}, que simula una huella de navegador real gestionando
    automáticamente los desafíos de JavaScript y cabeceras \textit{User-Agent}.
\end{itemize}
 
Para garantizar la robustez de las descargas, se implementó un algoritmo de
\textit{backoff} exponencial con varianza aleatoria, que evita que el servidor identifique
patrones de automatización y permite recuperarse de errores transitorios de red:
 
\begin{lstlisting}[language=Python,caption={Algoritmo de reintentos con backoff exponencial}]
delay_base = min(backoff_inicial * (2 ** (intento - 1)), backoff_max)
delay = delay_base + random.uniform(0, 1.5)
time.sleep(delay)
\end{lstlisting}
 
\subsubsection{8.1.2 Módulos de scraping}
 
La lógica de extracción se estructura en módulos especializados:
 
\begin{itemize}
    \item \textbf{\texttt{fbref.py}}: Parsea las tablas HTML de cada partido (summary,
    passing, defense, etc.) mediante BeautifulSoup. Combina múltiples tablas en memoria
    usando el nombre normalizado del jugador como clave de unión. Incluye lógica especializada
    para porteros, cuyas métricas (PSxG, Save\%, GA90) se extraen de tablas parcialmente
    distintas.
 
    \item \textbf{\texttt{roles.py}}: Extrae los rankings de eficiencia de FBRef
    (goles por disparo, pases al último tercio, intercepciones, etc.) desde nodos HTML
    ocultos en comentarios del DOM, usando una técnica de extracción específica:
 
\begin{lstlisting}[language=Python,caption={Extracción de datos ocultos en comentarios HTML}]
comments = soup.find_all(string=lambda text: isinstance(text, Comment))
for c in comments:
    if 'id="div_leaders"' in c:
        frag_soup = BeautifulSoup(c, "lxml")
\end{lstlisting}
 
    \item \textbf{\texttt{transfermarkt.py}}: Extrae la clasificación por jornada
    (PJ, PG, PE, PP, GF, GC, puntos) y datos de contexto competitivo.
 
    \item \textbf{\texttt{commons.py}}: Motor de normalización de texto compartido entre
    todos los módulos. Aplica normalización NFD para eliminar acentos y diacríticos,
    expresiones regulares para limpiar caracteres especiales, y detección de minutos
    concatenados al nombre del jugador (ej.: \texttt{90+4'}).
 
    \item \textbf{\texttt{popularDB.py}}: Orquestador final del pipeline. Valida la
    disponibilidad local de los HTMLs, crea las entidades base en la base de datos
    (temporadas, equipos, jornadas), ensambla las estadísticas por partido y jugador,
    aplica la normalización final y persiste los datos en PostgreSQL.
\end{itemize}
 
\subsection{8.2 Matching de entidades entre fuentes}
 
Uno de los mayores retos técnicos del proyecto fue la desambiguación de nombres entre
fuentes. Un mismo jugador puede aparecer como ``Lamine Yamal'', ``L. Yamal'' o ``Yamal''
según el portal. El proceso de \textit{matching} se articula en tres niveles.
 
\subsubsection{8.2.1 Normalización de texto}
 
Antes de cualquier comparación, el sistema aplica una normalización agresiva:
conversión a minúsculas, eliminación de tildes mediante NFD, eliminación de caracteres
especiales (guiones, puntos) y limpieza de minutos concatenados por FBRef al nombre
del jugador.
 
\subsubsection{8.2.2 Diccionario de alias}
 
Para casos donde la normalización es insuficiente (ej.: ``Pepelu'' vs. ``José Luis García
Vaya''), el sistema utiliza un diccionario de alias por temporada y equipo que intercepta
estos casos antes de la lógica probabilística.
 
\subsubsection{8.2.3 Fuzzy Matching contextual a tres niveles}
 
El \textit{matching} no compara contra toda la base de datos, sino que se acota al
contexto del partido (los jugadores de los dos equipos enfrentados), eliminando falsos
positivos masivos. La función \texttt{obten\_coincidencia\_nombre} aplica una estrategia
en cascada:
 
\begin{enumerate}
    \item \textbf{Regla de apellidos críticos}: Para apellidos muy comunes en el fútbol
    (García, González, Williams...) se exige coincidencia del nombre completo. Este nivel
    resuelve casos como los hermanos Iñaki y Nico Williams del Athletic Club.
 
    \item \textbf{Coincidencia por prefijo/sufijo}: Para nombres de un único token sin
    apellido crítico, el sistema detecta si coincide con el inicio o fin de algún candidato.
    Resuelve apodos y abreviaciones (``Araújo'' $\rightarrow$ ``Ronald Araújo'').
 
    \item \textbf{Fallback a WRatio}: Si los niveles anteriores no producen un resultado
    confiable, se recurre al algoritmo \textbf{WRatio} de la librería \texttt{rapidfuzz},
    que combina internamente \textit{Partial Ratio}, \textit{Token Sort Ratio} y escalado
    ponderado de pesos. Esto supera las limitaciones de la distancia de Levenshtein clásica,
    que rechazaría como válida la coincidencia ``Isaac Palazon'' frente a ``Isaac Palazon
    Camacho'' por la diferencia de caracteres.
\end{enumerate}
 
En caso de conflicto (el mismo jugador Fantasy propuesto para múltiples registros de FBRef),
el sistema aplica criterios de desempate por minutos jugados, puntuación Fantasy y
consistencia del apellido.
 
\subsection{8.3 Limpieza y saneamiento del dataset}
 
Una vez consolidado el histórico, el pipeline aplica las siguientes transformaciones de
saneamiento:
 
\begin{itemize}
    \item \textbf{Tratamiento de outliers}: Se identifican registros con puntuaciones
    anómalas superiores a 30 puntos en una jornada (habitualmente errores de las fuentes) y
    se sustituyen por la mediana o moda estadística del jugador, evitando que el modelo
    aprenda de eventos estadísticamente irrelevantes.
 
    \item \textbf{Filtrado de actividad}: Se eliminan registros de jugadores con menos de
    10 minutos disputados, ya que su aporte estadístico no es representativo y genera
    métricas por 90 minutos distorsionadas.
 
    \item \textbf{Filtrado por historial mínimo}: Se exige un mínimo de 5 partidos por
    jugador para que las variables de ventana deslizante tengan una base histórica
    estadísticamente fiable.
 
    \item \textbf{Sanitización de valores no válidos}: Los valores \texttt{NaN} e
    infinitos se corrigen mediante imputación por mediana o relleno controlado en función
    de la fase del procesamiento.
\end{itemize}
 
\subsection{8.4 Ingeniería de características}
 
El sistema no utiliza datos brutos para el entrenamiento, sino variables elaboradas que
capturan tendencias temporales y contexto competitivo. Es importante destacar que estas
variables derivadas no residen de forma persistente en la base de datos; se generan
dinámicamente durante el pipeline de entrenamiento a partir de los datos brutos, garantizando
que los cálculos se realicen siempre sobre el estado más actualizado y sin redundancia de
almacenamiento.
 
\subsubsection{8.4.1 Variables de serie temporal}
 
Para capturar el estado de forma de cada jugador, se generan métricas sobre ventanas
deslizantes con desplazamiento temporal (\texttt{shift(1)}), garantizando que cada fila
utilice únicamente información anterior al partido a predecir:
 
\begin{itemize}
    \item \textbf{Rolling Mean (media móvil simple)}: Promedio de las últimas $N$
    jornadas (ventanas de 3 y 5). Captura la estabilidad y el rendimiento sostenido.
    \item \textbf{EWMA (Exponential Weighted Moving Average)}: Media ponderada que
    prioriza los encuentros más recientes, permitiendo al modelo reaccionar con mayor
    sensibilidad ante cambios de forma, rol o minutaje.
\end{itemize}
 
$$EWMA_t = \alpha \cdot X_t + (1-\alpha) \cdot EWMA_{t-1}$$
 
Durante el desarrollo se evaluó también una ventana de 7 partidos, que fue descartada
por suavizar en exceso la señal y enmascarar cambios relevantes en el rendimiento.
 
\subsubsection{8.4.2 Enriquecimiento con roles de élite}
 
Los rankings de eficiencia extraídos de FBRef (jugadores en el top-10 de una métrica en
la liga) se convierten en variables numéricas mediante un factor de posición exponencial.
Un jugador en el puesto 1 de un ranking tiene un peso significativamente mayor que uno
en el puesto 10:
 
$$Score = \sum \left(Valor \times Peso \times \frac{1}{\sqrt{Posicion}}\right)$$
 
El módulo \texttt{role\_enricher.py} crea adicionalmente variables de interacción entre
la condición de élite y la intensidad del rol específico del jugador (ej.:
\texttt{elite\_paradas\_interact} para porteros).
 
\subsubsection{8.4.3 Integración de cuotas de mercado y contexto de partido}
 
Las cuotas de las casas de apuestas, transformadas en probabilidades implícitas, se
incorporan como señal del ``sentimiento del mercado'' previo al partido. Esta fue una
decisión de diseño determinante: en las fases iniciales del proyecto sin estas variables,
los resultados predictivos eran significativamente peores. Los datos se obtienen de
\href{https://www.football-data.co.uk/spainm.php}{football-data.co.uk}.
 
Variables derivadas de este bloque: probabilidad implícita de victoria (\texttt{p\_home},
\texttt{p\_away}), probabilidad de más de 2.5 goles (\texttt{p\_over25\_ewma5}), dificultad
del rival (\texttt{fixture\_difficulty\_home}, \texttt{fixture\_difficulty\_away}), y
goles encajados esperados por el rival (\texttt{opp\_gc\_ewma5}).
 
\subsubsection{8.4.4 Variables específicas por posición}
 
Para maximizar la precisión predictiva, se han diseñado conjuntos de variables específicos
según el rol del jugador. Las variables de ruido (estadísticas propias de otra posición)
se eliminan explícitamente para evitar que el modelo aprenda de señales sin valor
predictivo real para cada demarcación.
 
\paragraph{Porteros (PT/GK)}
 
Las variables de creación ofensiva (goles, xG, regates, conducciones, duelos) se eliminan
por ser ruido puro para esta posición. Las variables específicas generadas son:
 
\begin{itemize}
    \item \textbf{\texttt{cs\_probability}}: Probabilidad compuesta de portería a cero,
    calculada como media ponderada de cuatro señales: probabilidad implícita de las apuestas
    (peso 0.4), ratio inverso de goles encajados del rival (0.2), ratio inverso de tiros
    del rival (0.2) y dificultad del fixture (0.2).
    \item \textbf{\texttt{cs\_expected\_points}}: Puntos esperados por \textit{clean sheet}
    ($4 \times cs\_probability$).
    \item \textbf{\texttt{save\_per\_90\_ewma5}}: Media exponencial de paradas por cada
    90 minutos, capturando la actividad y eficacia del portero.
    \item \textbf{\texttt{psxg\_per\_90\_ewma5}}: Goles esperados tras el tiro (PSxG)
    por 90 minutos, indicador de la presión recibida.
    \item \textbf{\texttt{expected\_gk\_core\_points}}: Variable sintética que combina
    CS esperado, volumen de paradas y PSxG en un único indicador:
    $cs\_expected\_points + 0.4 \cdot save\_per\_90\_ewma5 - 0.3 \cdot psxg\_per\_90\_ewma5$.
\end{itemize}
 
\paragraph{Defensas (DF)}
 
Las métricas de portería (paradas, PSxG, goles en contra y sus derivadas temporales) se
eliminan como ruido. Las variables específicas generadas son:
 
\begin{itemize}
    \item \textbf{\texttt{tackles\_per\_90\_ewma5}}: Intensidad defensiva en entradas
    normalizada a 90 minutos, con suavizado exponencial.
    \item \textbf{\texttt{int\_per\_90\_ewma5}}: Capacidad de interceptación por 90
    minutos, con tendencia exponencial.
    \item \textbf{\texttt{clearances\_per\_90}}: Despejes normalizados a 90 minutos.
    \item \textbf{\texttt{defensive\_actions\_total}}: Agregado ponderado de acciones
    defensivas ($entradas + intercepciones + despejes \times 0.5$).
    \item \textbf{\texttt{def\_actions\_per\_90}} y \textbf{\texttt{def\_actions\_ewma5}}:
    Normalización e historial exponencial de las acciones defensivas totales.
    \item \textbf{\texttt{cs\_activity\_alignment}}: Correlación entre la probabilidad de
    portería a cero del equipo y la actividad defensiva propia del jugador.
    \item \textbf{\texttt{int\_momentum}}: Tendencia en el número de intercepciones
    (diferencia EWMA3 $-$ EWMA5).
    \item \textbf{\texttt{defensive\_context}}: Interacción entre la condición de local
    y la dificultad del fixture.
\end{itemize}
 
\paragraph{Mediocentros (MC)}
 
Las métricas exclusivas de portería (paradas, PSxG, \textit{clean sheet} esperado) se
eliminan como ruido. Las variables específicas generadas son:
 
\begin{itemize}
    \item \textbf{\texttt{pass\_attempts\_per\_90}}: Volumen de distribución de juego
    normalizado a 90 minutos.
    \item \textbf{\texttt{pass\_accuracy\_consistency}}: Estabilidad en el acierto de
    pase en las últimas jornadas (EWMA5 del porcentaje de pases completados).
    \item \textbf{\texttt{dribbles\_per\_90}}: Frecuencia de regates por 90 minutos.
    \item \textbf{\texttt{dribble\_intensity}}: Media móvil de conducciones totales.
    \item \textbf{\texttt{progressive\_action\_rate}}: Tasa de conducciones progresivas
    por 90 minutos.
    \item \textbf{\texttt{defensive\_contribution}}: EWMA5 de las acciones defensivas
    totales (entradas más intercepciones), capturando el rol recuperador del centrocampista.
    \item \textbf{\texttt{creative\_vs\_defensive}}: Ratio entre el promedio móvil de
    pases y el de entradas, que caracteriza el perfil ofensivo o defensivo del jugador.
    \item \textbf{\texttt{dribble\_confidence}}: EWMA5 del porcentaje de regates
    completados, indicador de confianza en la conducción.
    \item \textbf{\texttt{overall\_activity}}: Suma de pases, entradas, regates y
    conducciones, como índice global de implicación en el juego.
    \item \textbf{\texttt{playing\_time\_form\_combo}}: Producto del porcentaje de
    minutos jugados (EWMA5) por los puntos Fantasy (EWMA5), penalizando a jugadores
    con buen rendimiento pero escasa participación.
\end{itemize}
 
\paragraph{Delanteros (DT)}
 
Las métricas de portería y las acciones defensivas (despejes, bloqueos, duelos aéreos,
\textit{clearances} temporales) se eliminan como ruido. Las variables específicas
generadas son:
 
\begin{itemize}
    \item \textbf{\texttt{xg\_overperformance}}: Diferencia entre goles marcados y xG
    esperado en los últimos 5 partidos, detectando estados de gracia o mala racha ante
    el gol.
    \item \textbf{\texttt{shot\_conversion\_ewma5}}: Eficiencia de cara a puerta
    (goles por tiro), con suavizado exponencial para capturar la tendencia reciente.
    \item \textbf{\texttt{offensive\_pressure\_score}}: Variable compuesta que pondera
    el volumen de tiros y el xG generado ($shots\_roll5 \times 0.4 + xg\_roll5 \times 0.6$).
    \item \textbf{\texttt{scoring\_stability}}: Inversa de la diferencia absoluta entre
    la media móvil de goles (roll5 y ewma5), midiendo la consistencia goleadora.
    \item \textbf{\texttt{xg\_per\_minute\_ewma}}: xG generado por minuto disponible,
    que penaliza a jugadores con buenos números pero escasa participación.
    \item \textbf{\texttt{goals\_per\_minute\_ewma}}: Goles por minuto disponible.
    \item \textbf{\texttt{scoring\_momentum}}: Tendencia de racha goleadora calculada
    como la diferencia EWMA3 $-$ EWMA5 de goles.
    \item \textbf{\texttt{xg\_momentum}}: Tendencia del xG generado (EWMA3 $-$ EWMA5).
    \item \textbf{\texttt{pf\_consistency\_fw}}: Inversa de la volatilidad de los puntos
    Fantasy en los últimos 5 partidos, midiendo la fiabilidad del delantero.
\end{itemize}
 
\subsubsection{8.4.5 Variables comunes eliminadas por ruido}
 
Con independencia de la posición, se eliminan sistemáticamente aquellas variables que
los modelos penalizan con importancia nula o que introducen colinealidad sin valor
predictivo adicional: \texttt{duels\_won\_pct\_ewma5}, \texttt{duels\_won\_roll5},
\texttt{duels\_won\_ewma5}, \texttt{elite\_entradas\_interact} y
\texttt{ratio\_roles\_criticos}.
 
\subsection{8.5 Prevención del data leakage}
 
Para garantizar que el sistema sea aplicable en el mundo real, se implementan dos
salvaguardas fundamentales:
 
\begin{itemize}
    \item \textbf{Desplazamiento temporal (Shift)}: Todas las estadísticas de rendimiento
    se desplazan una jornada hacia atrás (\texttt{shift(1)}) por jugador. Para predecir
    la jornada $N$, el modelo solo tiene acceso a variables calculadas hasta la jornada
    $N-1$.
 
    \item \textbf{Validación temporal estricta}: En lugar de una validación aleatoria
    (K-Fold), se utiliza un corte temporal: el 80\% de los datos más antiguos se usan
    para entrenamiento y el 20\% más reciente para test. Esto garantiza que el modelo
    nunca aprende del futuro.
\end{itemize}
 
\subsection{8.6 Estructura del dataset final}
 
Tras el proceso de ETL y enriquecimiento, el dataset final presenta una estructura robusta
con un total de \textbf{29.222 filas} (observaciones jugador/partido) y \textbf{68 columnas}.
Las variables se agrupan en los siguientes bloques:
 
\begin{table}[H]
\centering
\renewcommand{\arraystretch}{1.3}
\begin{tabular}{|p{4.5cm}|p{9cm}|}
\hline
\textbf{Bloque} & \textbf{Variables representativas} \\
\hline
Identificación y Bio &
\texttt{player, nacionalidad, edad, posicion, dorsal} \\
\hline
Contexto de Encuentro &
\texttt{equipo\_propio, equipo\_rival, is\_home, temporada, jornada, fecha\_partido} \\
\hline
Rendimiento ofensivo &
\texttt{gol\_partido, asist\_partido, xg\_partido, xag, tiros, tiro\_puerta\_partido,
regates\_completados, conducciones\_progresivas, pases\_totales,
pases\_completados\_pct} \\
\hline
Métricas defensivas y portería &
\texttt{entradas, intercepciones, despejes, bloqueos, duelos\_ganados,
porcentaje\_paradas, psxg, goles\_en\_contra} \\
\hline
Contexto competitivo &
\texttt{pos\_clasificacion, pj, pg, pe, pp, gf, gc, pts, racha5partidos,
opp\_gc\_ewma5, opp\_shots\_ewma5, fixture\_difficulty\_home/away} \\
\hline
Cuotas de mercado &
\texttt{p\_home, p\_away, p\_over25, p\_over25\_ewma5} \\
\hline
Disponibilidad y forma &
\texttt{min\_partido, minutes\_pct\_roll3/5, minutes\_pct\_ewma3/5,
pf\_roll3/5, pf\_ewma3/5} \\
\hline
Variable objetivo &
\texttt{target\_pf\_next} \\
\hline
\end{tabular}
\caption{Bloques de variables del dataset final.}
\label{tab:dataset-estructura}
\end{table}
 
\noindent Las variables de serie temporal (rolling y EWMA) y las features específicas por posición
descritas en el apartado 8.4.4 se generan dinámicamente sobre este conjunto base durante
la fase de entrenamiento, y no se almacenan de forma persistente en la base de datos.
